\section{Background}
\label{sec:background}

\subsection{Field-Level Conformance}

Protocol conformance asks whether an implementation satisfies the behavioral
requirements stated by a standard. In natural-language standards, many
requirements are attached to protocol fields: message fields, parameters,
lengths, flags, and valid field combinations. We call this scope
\emph{field-level conformance}.

Field-level conformance requires two links: mapping the standard's field and
condition to implementation objects, and checking whether the relevant code path
enforces the required behavior. This scope is narrower than full protocol
verification, but it covers many practical deviations because parsing,
validation, negotiation, and error handling are often organized around fields.

\subsection{Normative Modalities}

Normative modality determines whether a statement is treated as a reportable
obligation. \tool{} extracts \emph{MUST} and \emph{MUST NOT} requirements as
primary conformance rules, tracks \emph{SHOULD} requirements as lower-severity
deviations, and ignores \emph{MAY} statements because they describe optional
behavior. The modality is stored with each rule and used when ranking reports.

% \subsection{A Running Example}

% Consider a simplified QUIC transport rule: if an endpoint receives a transport
% parameter that violates a specified bound, it must terminate the connection with
% a protocol error. The standard may express this requirement across a field
% definition, a paragraph about transport parameters, and an error-handling
% section. The implementation may parse the parameter in one file, store it in a
% connection configuration structure, validate it in a helper, and trigger the
% error through a callback. A direct text-to-code query such as ``does the code
% enforce the transport parameter bound?'' is under-specified: the relevant
% variable names, call chain, and build configuration are all missing.

% \tool{} handles the example by first extracting a rule such as
% \[
% \begin{aligned}
% C &= \textit{parameter received},\\
% f &= \textit{max\_udp\_payload\_size},\\
% A &= \textit{reject values below the lower bound}.
% \end{aligned}
% \]
% It then searches for identifiers and code regions related to the field,
% including aliases such as \texttt{max\_udp\_payload},
% \texttt{udp\_payload\_size}, or structure members embedded in parser state.
% If the first retrieved context only shows assignment but not validation, the
% reasoning agent emits \emph{insufficient evidence} and requests callers or
% validation helpers. Only after the evidence covers the relevant path does the
% system report a consistency judgment. For a suspected deviation, the validation-test agent
% attempts to construct a packet or configuration that exercises the suspicious
% path and observes whether the implementation accepts or rejects it.

% \subsection{Threats to Naive LLM Checking}

% LLMs are useful because they can relate prose to code-level intent, but
% conformance checking punishes ungrounded flexibility. First, field names are
% ambiguous: a standard may use a short name in one paragraph and a fully
% qualified message path in another. Second, code identifiers are indirect:
% \texttt{len}, \texttt{payload\_len}, and \texttt{param->value} may refer to
% different protocol objects depending on scope. Third, evidence is often
% negative: a missing check is meaningful only after the checker has searched
% the right parser, validation helper, and caller constraints. Finally, execution
% semantics may depend on macros, compile-time options, and runtime configuration.
% \tool{} is designed around these threats. It constrains field references,
% retrieves code structurally, records uncertainty, and validates high-risk
% claims dynamically when possible.
